% GENERATED FILE. Edit canonical lessons, then run npm run generate:books.
% canonical-source-hash: fnv1a64:420f9620bc05fac2
% canonical-lessons: LA-C534-osculum, LA-C534-basium, LA-C534-basio, LA-C534-suavium, LA-C534-amplexus

\chapter{Kisses and hugs}
\label{ch:la-kisses-and-hugs}
\hlchaptermodality{534}

\begin{chapteropening}
\textbf{By the end of this chapter:} \emph{I can talk about kisses and hugs.}

Complete the last lesson of chapter 534: I can talk about kisses and hugs.
\end{chapteropening}

\section[ōsculum, ōsculī]{ōsculum, ōsculī — a kiss}

\label{lesson:LA-C534-osculum}



\begin{quote}
Before the new one: say the Latin for to hold dear, then the Latin for lovable.
\end{quote}

\subsection*{You'll want to know: ōsculum, ōsculī}
\textbf{ōsculum, ōsculī} — \textquotedblleft{}a kiss\textquotedblright{}.

\begin{quote}
Māter fīliō ōsculum dat.
\end{quote}

\begin{quote}
\emph{The mother gives her son a kiss.}
\end{quote}

Neuter, a diminutive of \emph{ōs} — mouth: a little mouth. \emph{Ōsculor}, already known, is to kiss.

\begin{cousinweb}
\textbf{ōsculum} is a diminutive of \emph{ōs} \textquotedblleft{}mouth\textquotedblright{} (Lewis \& Short; EtymDB). English \emph{osculate} comes from the verb \emph{ōsculor}.
\end{cousinweb}

\begin{grammarlens}[title={What you've built}]
One more word for telling, explaining and describing at B1.
\end{grammarlens}

\subsection*{Your turn}
\begin{itemize}
\raggedright
  \item \emph{Say it:} \emph{ōsculum, ōsculī}
  \item \emph{Say it:} \emph{ōsculum, ōsculī}, once more
  \item \emph{Say it:} a sentence of your own with \emph{ōsculum, ōsculī}
  \item \emph{Recall:} say the Latin for to hold dear, then the Latin for lovable, then say \emph{ōsculum, ōsculī} again
\end{itemize}

\begin{tcolorbox}[breakable,colback=teal!4,colframe=teal!35!black,title={Before you move on}]
What does ōsculum, ōsculī mean? (\textquotedblleft{}A kiss\textquotedblright{}.) Say it once more.
\end{tcolorbox}
\section[bāsium, bāsiī]{bāsium, bāsiī — a loving kiss}

\label{lesson:LA-C534-basium}



\begin{quote}
Before the new one: say the Latin for lovable, then the Latin for a kiss.
\end{quote}

\subsection*{You'll want to know: bāsium, bāsiī}
\textbf{bāsium, bāsiī} — \textquotedblleft{}a loving kiss\textquotedblright{}.

\begin{quote}
Fīliola patrī bāsia mīlle dat.
\end{quote}

\begin{quote}
\emph{The little girl gives her father a thousand kisses.}
\end{quote}

Neuter, second declension. A warm word of love poetry. Spanish \emph{beso} and Italian \emph{bacio} come from it.

\begin{cousinweb}
\textbf{bāsium} is a word of uncertain origin; Lewis \& Short compares English \emph{buss}, and EtymDB gives a Proto-Indo-European root \emph{bu-}, but it is not settled. Spanish \emph{beso} and Italian \emph{bacio} come from it.
\end{cousinweb}

\begin{grammarlens}[title={What you've built}]
One more word for telling, explaining and describing at B1.
\end{grammarlens}

\subsection*{Your turn}
\begin{itemize}
\raggedright
  \item \emph{Say it:} \emph{bāsium, bāsiī}
  \item \emph{Say it:} \emph{bāsium, bāsiī}, once more
  \item \emph{Say it:} a sentence of your own with \emph{bāsium, bāsiī}
  \item \emph{Recall:} say the Latin for lovable, then the Latin for a kiss, then say \emph{bāsium, bāsiī} again
\end{itemize}

\begin{tcolorbox}[breakable,colback=teal!4,colframe=teal!35!black,title={Before you move on}]
What does bāsium, bāsiī mean? (\textquotedblleft{}A loving kiss\textquotedblright{}.) Say it once more.
\end{tcolorbox}
\section[bāsiō, bāsiāre]{bāsiō, bāsiāre — to kiss}

\label{lesson:LA-C534-basio}



\begin{quote}
Before the new one: say the Latin for a kiss, then the Latin for a loving kiss.
\end{quote}

\subsection*{You'll want to know: bāsiō, bāsiāre}
\textbf{bāsiō, bāsiāre} — \textquotedblleft{}to kiss\textquotedblright{}.

\begin{quote}
Avia nepōtēs bāsiat.
\end{quote}

\begin{quote}
\emph{Grandmother kisses her grandchildren.}
\end{quote}

\emph{bāsiō, bāsiāre, bāsiāvī, bāsiātum}, from \emph{bāsium}.

\begin{cousinweb}
\textbf{bāsiō} is built on \emph{bāsium} \textquotedblleft{}a kiss\textquotedblright{} (Lewis \& Short; EtymDB). French \emph{baiser} comes from it.
\end{cousinweb}

\begin{grammarlens}[title={What you've built}]
One more word for telling, explaining and describing at B1.
\end{grammarlens}

\subsection*{Your turn}
\begin{itemize}
\raggedright
  \item \emph{Say it:} \emph{bāsiō, bāsiāre}
  \item \emph{Say it:} \emph{bāsiō, bāsiāre}, once more
  \item \emph{Say it:} a sentence of your own with \emph{bāsiō, bāsiāre}
  \item \emph{Recall:} say the Latin for a kiss, then the Latin for a loving kiss, then say \emph{bāsiō, bāsiāre} again
\end{itemize}

\begin{tcolorbox}[breakable,colback=teal!4,colframe=teal!35!black,title={Before you move on}]
What does bāsiō, bāsiāre mean? (\textquotedblleft{}To kiss\textquotedblright{}.) Say it once more.
\end{tcolorbox}
\section[suāvium, suāviī]{suāvium, suāviī — a kiss on the lips}

\label{lesson:LA-C534-suavium}



\begin{quote}
Before the new one: say the Latin for a loving kiss, then the Latin for to kiss.
\end{quote}

\subsection*{You'll want to know: suāvium, suāviī}
\textbf{suāvium, suāviī} — \textquotedblleft{}a kiss on the lips\textquotedblright{}.

\begin{quote}
Puella iuvenī suāvium dat.
\end{quote}

\begin{quote}
\emph{The girl gives the young man a kiss on the lips.}
\end{quote}

Neuter, second declension; also spelt \emph{sāvium}. A word of comedy and love letters.

\begin{cousinweb}
\textbf{suāvium} is built on \emph{suāvis} \textquotedblleft{}sweet\textquotedblright{} (Lewis \& Short).
\end{cousinweb}

\begin{grammarlens}[title={What you've built}]
One more word for telling, explaining and describing at B1.
\end{grammarlens}

\subsection*{Your turn}
\begin{itemize}
\raggedright
  \item \emph{Say it:} \emph{suāvium, suāviī}
  \item \emph{Say it:} \emph{suāvium, suāviī}, once more
  \item \emph{Say it:} a sentence of your own with \emph{suāvium, suāviī}
  \item \emph{Recall:} say the Latin for a loving kiss, then the Latin for to kiss, then say \emph{suāvium, suāviī} again
\end{itemize}

\begin{tcolorbox}[breakable,colback=teal!4,colframe=teal!35!black,title={Before you move on}]
What does suāvium, suāviī mean? (\textquotedblleft{}A kiss on the lips\textquotedblright{}.) Say it once more.
\end{tcolorbox}
\section[amplexus, amplexūs]{amplexus, amplexūs — an embrace}

\label{lesson:LA-C534-amplexus}



\begin{quote}
Before the new one: say the Latin for to kiss, then the Latin for a kiss on the lips.
\end{quote}

\subsection*{You'll want to know: amplexus, amplexūs}
\textbf{amplexus, amplexūs} — \textquotedblleft{}an embrace\textquotedblright{}.

\begin{quote}
Māter fīliam in amplexū tenet.
\end{quote}

\begin{quote}
\emph{The mother holds her daughter in an embrace.}
\end{quote}

Masculine, fourth declension: arms wound around someone. In poetry also the coils of a snake.

\begin{cousinweb}
\textbf{amplexus} is the action noun of \emph{amplector} \textquotedblleft{}to embrace\textquotedblright{} (Lewis \& Short).
\end{cousinweb}

\begin{grammarlens}[title={What you've built}]
One more word for telling, explaining and describing at B1.
\end{grammarlens}

\subsection*{Your turn}
\begin{itemize}
\raggedright
  \item \emph{Say it:} \emph{amplexus, amplexūs}
  \item \emph{Say it:} \emph{amplexus, amplexūs}, once more
  \item \emph{Say it:} a sentence of your own with \emph{amplexus, amplexūs}
  \item \emph{Recall:} say the Latin for to kiss, then the Latin for a kiss on the lips, then say \emph{amplexus, amplexūs} again
\end{itemize}

\begin{tcolorbox}[breakable,colback=teal!4,colframe=teal!35!black,title={Before you move on}]
What does amplexus, amplexūs mean? (\textquotedblleft{}An embrace\textquotedblright{}.) Say it once more.
\end{tcolorbox}
